\section{Moving geometry: why pruning can be premature}
In a feature-learning network, the basis itself moves. For simple eigenvalues,
\[
\dot\mu_j=\ip{u_j}{\dot L_tu_j},
\]
and under the gauge $\ip{u_j}{\dot u_j}=0$,
\[
\dot u_j
=\sum_{m\ne j}
\frac{\ip{u_m}{\dot L_tu_j}}{\mu_j-\mu_m}u_m.
\]
The residual coefficients satisfy
\[
\boxed{
\dot a_j
=-\mu_ja_j
+\sum_{m\ne j}
\frac{\ip{u_m}{\dot L_tu_j}}{\mu_j-\mu_m}a_m.}
\]
The first term is coefficient fitting in the current basis; the second is target transport caused by feature learning. A currently weak mode can therefore become important because its eigenvalue grows, its eigenspace rotates toward residual target mass, or both.

This fact motivates a \emph{persistence} requirement: deletion should be based on a finite-horizon upper bound, not a single checkpoint.

\section{Curvature controls how fast prediction geometry can move}
The network does not move its prediction operator arbitrarily fast. The exact gradient-dissipation identity already limits the distance that the geometry can travel when curvature is controlled.

For the second Fr\'echet derivative define
\[
\norm{D^2F(\theta)}
:=\sup_{\norm v=\norm w=1}\norm{D^2F(\theta)[v,w]}_{\Hh}.
\]

\begin{theorem}[Curvature-to-spectral-motion bound]\label{thm:curvature-motion}
Along square-loss gradient flow, assume on $[t,t+H]$ that
\[
\norm{J_s}_{\op}\le B_J,
\qquad
\norm{D^2F(\theta_s)}\le B_H.
\]
Then
\[
\boxed{
\norm{\dot L_s}_{\op}
\le 2B_HB_J\norm{J_s^\ast e_s}.}
\]
Consequently the operator path length satisfies
\[
\boxed{
\int_t^{t+H}\norm{\dot L_s}_{\op}\,ds
\le
2B_HB_J\sqrt{H\,[\R(\theta_t)-\R(\theta_{t+H})]}
\le
2B_HB_J\sqrt{H\R(\theta_t)}.}
\]
If $P_I(s)$ is the Riesz projector of a spectral cluster separated from the rest of $L_s$ by a gap at least $\Delta>0$, then
\[
\boxed{
\int_t^{t+H}\norm{\dot P_I(s)}_{\op}\,ds
\le
\frac{4B_HB_J}{\Delta}
\sqrt{H\,[\R(\theta_t)-\R(\theta_{t+H})]}.}
\]
\end{theorem}

\begin{proof}
Differentiating $J_s=DF(\theta_s)$ along the trajectory gives
\[
\dot J_s[v]=D^2F(\theta_s)[\dot\theta_s,v],
\]
so
\[
\norm{\dot J_s}_{\op}\le B_H\norm{\dot\theta_s}
=B_H\norm{J_s^\ast e_s}.
\]
Since $L_s=J_sJ_s^\ast$,
\[
\dot L_s=\dot J_sJ_s^\ast+J_s\dot J_s^\ast,
\]
and hence
\[
\norm{\dot L_s}_{\op}
\le2\norm{J_s}_{\op}\norm{\dot J_s}_{\op}
\le2B_HB_J\norm{J_s^\ast e_s}.
\]
By Theorem~\ref{thm:grad-energy},
\[
\int_t^{t+H}\norm{J_s^\ast e_s}^2ds
=\R(\theta_t)-\R(\theta_{t+H}).
\]
Cauchy--Schwarz proves the operator path-length bound. Finally, the standard differentiable Riesz-projector bound
$\norm{\dot P_I}_{\op}\le2\norm{\dot L}_{\op}/\Delta$
gives the cluster path-length claim.
\end{proof}

\begin{remark}
This result supplies the missing dynamical input to a feature-value stopping argument. A large unresolved target component is not enough to justify continued full-network training: the network must also have enough remaining curvature-driven subspace motion to reach it. Conversely, early in training a large geometry path budget is a mathematical reason not to prune merely because current eigenvalues are small.
\end{remark}

\section{Finite-horizon stability of structural contraction}
We now compare nonlinear full training with a physically reduced flow. Define
\[
g(\theta)=J_\theta^\ast e_\theta=-\nabla\R(\theta).
\]
Let $S$ be a fixed orthogonal projector representing the retained structural coordinates after contraction. We allow an instantaneous approximation error at the contraction checkpoint, because a reduced architecture need not be an exact embedding of the dense state. Starting at time $t$, write
\[
\dot\theta_s=g(\theta_s),
\qquad
\dot{\widetilde\theta}_s=Sg(\widetilde\theta_s),
\qquad
\norm{\theta_t-\widetilde\theta_t}=d_0.
\]
The reduced parameters may be embedded in the original parameter space by setting deleted structural coordinates inactive.

\begin{theorem}[Safe finite-horizon structural-flow bound]\label{thm:gronwall}
Assume on a tube containing both trajectories for $s\in[t,t+H]$ that:
\begin{enumerate}[label=(\roman*)]
\item $g$ is $L_g$-Lipschitz;
\item $\norm{J_\theta}_{\op}\le B_J$;
\item both residual norms are at most $B_e$;
\item the omitted dense-path gradient obeys
$\norm{(I-S)g(\theta_s)}\le\varepsilon(s)$.
\end{enumerate}
Then
\[
\boxed{
\norm{\theta_{t+H}-\widetilde\theta_{t+H}}
\le
e^{L_gH}d_0
+\int_t^{t+H}e^{L_g(t+H-s)}\varepsilon(s)\,ds,}
\]
\[
\boxed{
\norm{F(\theta_{t+H})-F(\widetilde\theta_{t+H})}
\le
B_J\left[
e^{L_gH}d_0
+\int_t^{t+H}e^{L_g(t+H-s)}\varepsilon(s)\,ds
\right],}
\]
and
\[
\boxed{
|\R(\theta_{t+H})-\R(\widetilde\theta_{t+H})|
\le
B_eB_J\left[
e^{L_gH}d_0
+\int_t^{t+H}e^{L_g(t+H-s)}\varepsilon(s)\,ds
\right].}
\]
If $d_0=0$ and $\varepsilon(s)\le\bar\varepsilon$, the dynamic term is
$\bar\varepsilon(e^{L_gH}-1)/L_g$, with continuous limit $H\bar\varepsilon$ when $L_g=0$.
\end{theorem}

\begin{proof}
Set $z_s=\theta_s-\widetilde\theta_s$. Since $S$ is a contraction,
\[
\dot z_s
=(I-S)g(\theta_s)+S[g(\theta_s)-g(\widetilde\theta_s)],
\]
so
\[
\frac{d}{ds}\norm{z_s}
\le\varepsilon(s)+L_g\norm{z_s}
\]
almost everywhere. Gr\"onwall's inequality gives the parameter bound including the initial discrepancy $d_0$. The Jacobian bound makes $F$ $B_J$-Lipschitz on the tube. Finally,
\[
|\R(\theta)-\R(\tilde\theta)|
\le\frac12(\norm{e_\theta}+\norm{e_{\tilde\theta}})
\norm{e_\theta-e_{\tilde\theta}}
\le B_e\norm{F(\theta)-F(\tilde\theta)}.
\]
\end{proof}

This formulation is useful operationally because the omitted-gradient certificate may be estimated or bounded along the dense reference trajectory before an irreversible contraction is executed.

\section{A spectral persistence certificate}
Let $P_D(t)$ be the Riesz projector onto an isolated positive spectral cluster $D$ of $L_t$, separated from the rest of the spectrum by a gap $\Delta_t\ge\Delta>0$. Define the cluster gradient energy
\[
E_D(t)=\norm{J_t^\ast P_D(t)e_t}^2
=\ip{e_t}{L_tP_D(t)e_t}.
\]
For a simple spectrum this is $\sum_{j\in D}\mu_j a_j^2$.

\begin{theorem}[Upper growth bound for discarded-cluster gradient energy]\label{thm:persistence}
Suppose $L_t$ is differentiable, $P_D(t)$ remains isolated by gap $\Delta$, and square-loss gradient flow holds. Then
\[
\dot E_D(t)
\le
\norm{e_t}^2\norm{\dot L_t}_{\op}
\left(1+\frac{2\norm{L_t}_{\op}}{\Delta}\right).
\]
If on $[t,t+H]$
\[
\norm{e_s}\le B_e,
\qquad
\norm{L_s}_{\op}\le B_L,
\]
then for $0\le h\le H$,
\[
\boxed{
E_D(t+h)
\le E_D(t)
+B_e^2\left(1+\frac{2B_L}{\Delta}\right)
\int_t^{t+h}\norm{\dot L_s}_{\op}\,ds.}
\]
\end{theorem}

\begin{proof}
Because $P_D$ is a spectral projector, it commutes with $L$. Put $A=LP_D$. Then
\[
E_D=\ip e{Ae},
\quad
\dot E_D=2\ip{\dot e}{Ae}+\ip e{\dot A e}.
\]
Using $\dot e=-Le$ and $LP_D=P_DL$,
\[
2\ip{\dot e}{Ae}
=-2\ip{Le}{LP_De}
=-2\norm{LP_De}^2\le0.
\]
Also $\dot A=\dot L P_D+L\dot P_D$, so
\[
\dot E_D\le\norm e^2
\left(\norm{\dot L}_{\op}+\norm L_{\op}\norm{\dot P_D}_{\op}\right).
\]
The Riesz-projector/Davis--Kahan derivative bound gives
$\norm{\dot P_D}_{\op}\le2\norm{\dot L}_{\op}/\Delta$. Integrating proves the claim.
\end{proof}

\begin{corollary}[Curvature-based finite-horizon omitted-energy envelope]\label{cor:energy-envelope}
Under Theorems~\ref{thm:curvature-motion} and~\ref{thm:persistence}, define
\[
C_D=B_e^2\left(1+\frac{2B_L}{\Delta}\right),
\qquad
V_L(t,H)=2B_HB_J\sqrt{H[\R(\theta_t)-\R(\theta_{t+H})]}.
\]
Then
\[
E_D(t+h)\le E_D(t)+C_DV_L(t,H),
\qquad 0\le h\le H,
\]
and
\[
\boxed{
\int_t^{t+H}E_D(s)\,ds
\le H\,[E_D(t)+C_DV_L(t,H)].}
\]
In particular,
\[
\int_t^{t+H}\sqrt{E_D(s)}\,ds
\le
H\sqrt{E_D(t)+C_DV_L(t,H)}.
\]
\end{corollary}

The gap is essential. Inside an unresolved random-matrix bulk, individual empirical eigenvectors are unstable; deletion decisions should be made at the cluster level until a stable subspace detaches.

\section{From ideal spectral reduction to actual parameter deletion}
A spectral projector is not automatically a hardware-friendly parameter mask. This is the structural-realizability gap.

Let $Q_K$ be an ideal kept parameter projector and let $S$ be a structural projector induced by a realizable architecture reduction (for example, retained neurons or channels).

\begin{theorem}[Structural transfer]\label{thm:structural-transfer}
At any checkpoint,
\[
\boxed{
\norm{(I-S)g}
\le
\norm{(I-Q_K)g}
+\norm{(I-S)Q_Kg}.}
\]
In particular, if the ideal dropped sector has energy
$\norm{(I-Q_K)g}^2\le E_D$ and the structural approximation satisfies
$\norm{(I-S)Q_Kg}\le\delta$, then
\[
\norm{(I-S)g}\le\sqrt{E_D}+\delta.
\]
\end{theorem}

\begin{proof}
Write
$(I-S)g=(I-S)(I-Q_K)g+(I-S)Q_Kg$ and use that an orthogonal projector has operator norm one.
\end{proof}

The pointwise theorem is enough for a checkpoint. For future safety, however, a structural mask chosen at time $t$ must continue to approximate the moving kept spectral subspace.

\begin{theorem}[Persistence of a fixed structural realization]\label{thm:structural-persistence}
Let $G_s=J_s^\ast J_s$ be the parameter Gram operator. Suppose its ideal kept cluster has spectral projector $Q_K(s)$ and remains separated by a gap at least $\Delta_G>0$ on $[t,t+H]$. Let $S$ be a fixed structural projector selected at time $t$, and define
\[
\delta_0=\norm{(I-S)Q_K(t)}_{\op}.
\]
Under the curvature assumptions of Theorem~\ref{thm:curvature-motion}, for $0\le h\le H$,
\[
\boxed{
\norm{(I-S)Q_K(t+h)}_{\op}
\le
\delta_0
+\frac{4B_HB_J}{\Delta_G}
\sqrt{h\,[\R(\theta_t)-\R(\theta_{t+h})]}.}
\]
Consequently, with $Q_D(s)=I-Q_K(s)$,
\[
\boxed{
\norm{(I-S)g(\theta_s)}
\le
\norm{Q_D(s)g(\theta_s)}
+\delta_S(s)\norm{g(\theta_s)},}
\]
where $\delta_S(s)$ is bounded by the preceding display.
\end{theorem}

\begin{proof}
As for $L_s$, differentiating $G_s=J_s^\ast J_s$ gives
\[
\norm{\dot G_s}_{\op}
\le2B_HB_J\norm{J_s^\ast e_s}.
\]
The Riesz-projector derivative bound in parameter space yields
\[
\norm{Q_K(t+h)-Q_K(t)}_{\op}
\le\frac{2}{\Delta_G}\int_t^{t+h}\norm{\dot G_s}_{\op}ds.
\]
Apply Cauchy--Schwarz and the exact loss-dissipation identity to obtain the first claim. Finally,
\[
(I-S)g=(I-S)Q_Dg+(I-S)Q_Kg
\]
and the projector contraction bound gives the second claim.
\end{proof}

\begin{corollary}[Modular safe-contraction certificate]\label{cor:master-certificate}
A sufficient finite-horizon certificate for a structural contraction is obtained by the following chain:
\[
\begin{gathered}
\text{curvature + current loss}
\Rightarrow
\text{spectral path budget}
\Rightarrow
\text{discarded-energy persistence}\\
\Downarrow\\
\text{structural omitted-gradient budget}
\Rightarrow
\text{terminal risk bound}
\end{gathered}
\]
Specifically, Theorem~\ref{thm:curvature-motion} bounds future spectral motion, Theorem~\ref{thm:persistence} bounds the ideal discarded energy, Theorem~\ref{thm:structural-persistence} controls a fixed structural realization, and Theorem~\ref{thm:gronwall} converts the resulting $\varepsilon(s)$ together with any initial contraction distortion $d_0$ into an explicit terminal prediction/risk error.
\end{corollary}

This forces an honest separation between two scientific claims: prediction geometry may say that a sector is expendable, while a particular neuron/channel realization may still approximate that sector poorly. The structural gap is not a nuisance to hide; it is part of the theorem.
